% from ega3-1-fr.tex
\oldpage[III]{151}

Dans tous les cas d'application connus, $A$ est en fait un anneau adique local
noethérien, donc les $Y_n$ sont des spectres d'anneaux artiniens locaux, et les
résultats de ce paragraphe et des précédents réduisent dans une large mesure la
géométrie algébrique sur un anneau adique local noethérien, à la géométrie algébrique
sur des anneaux locaux artiniens.
\end{env}

\begin{corollary}[5.1.8]
\label{III.5.1.8-fr}
Sous les conditions de (\hyperref[III.5.1.4-fr]{5.1.4}), l'application
$Z\mapsto\widehat Z=Z_{/(Z\cap X')}$ est une bijection de l'ensemble des
sous-préschémas fermés $Z$ de $X$, propres sur $Y$, sur l'ensemble des
sous-préschémas formels fermés de $\widehat X$, propres sur $\widehat Y$.
\end{corollary}

\begin{proof}
En effet, un sous-préschéma formel fermé de $\widehat X$ est de la forme
$(T,(\mathcal O_{\widehat X}/\mathcal A)|T)$, où $\mathcal A$ est un Idéal
cohérent de $\mathcal O_{\widehat X}$ (I, 10.14.2)~; si $T$ est propre sur
$\widehat Y$, il résulte de (\hyperref[III.5.1.4-fr]{5.1.4}) que
$\mathcal O_{\widehat X}/\mathcal A$ est isomorphe à un
$\mathcal O_{\widehat X}$-Module de la forme $\widehat{\mathcal F}$, où
$\mathcal F$ est un $\mathcal O_X$-Module cohérent de support propre sur $Y$~; en
outre, il résulte de (\hyperref[III.5.1.3-fr]{5.1.3}) que l'homomorphisme canonique
$\mathcal O_{\widehat X}\to\mathcal O_{\widehat X}/\mathcal A$ est de la forme
$\widehat u$, où $u:\mathcal O_X\to\mathcal F$ est un homomorphisme surjectif de
$\mathcal O_X$-Modules. Donc $\mathcal F$ est de la forme
$\mathcal O_X/\mathcal N$, où $\mathcal N$ est un Idéal cohérent de
$\mathcal O_X$, et $\mathcal A=\widehat{\mathcal N}$ (I, 10.8.8), d'où la
conclusion (I, 10.14.7).
\end{proof}

\subsection{Démonstration du théorème d'existence : cas projectif et quasi-projectif}
\label{subsection:III.5.2-fr}

\begin{env}[5.2.1]
\label{III.5.2.1-fr}
Sous les conditions de (\hyperref[III.5.1.4-fr]{5.1.4}), nous dirons provisoirement
qu'un $\mathcal O_{\widehat X}$-Module cohérent est algébrisable s'il est isomorphe
à un complété $\widehat{\mathcal F}$ d'un $\mathcal O_X$-Module cohérent
$\mathcal F$ de support propre sur $Y$.
\end{env}

\begin{lemma}[5.2.2]
\label{III.5.2.2-fr}
Soient $\mathcal F'$, $\mathcal G'$ deux $\mathcal O_{\widehat X}$-Modules
algébrisables. Pour tout homomorphisme $u:\mathcal F'\to\mathcal G'$,
$\operatorname{Ker}(u)$, $\operatorname{Im}(u)$ et $\operatorname{Coker}(u)$ sont
algébrisables.
\end{lemma}

\begin{proof}
En effet, si $\mathcal F'=\widehat{\mathcal F}$,
$\mathcal G'=\widehat{\mathcal G}$, où $\mathcal F$ et $\mathcal G$ sont des
$\mathcal O_X$-Modules cohérents de supports propres sur $Y$, on a
$u=\widehat v$, où $v:\mathcal F\to\mathcal G$ est un homomorphisme
(\hyperref[III.5.1.3-fr]{5.1.3}). En vertu de l'exactitude du foncteur
$\mathcal F\mapsto\widehat{\mathcal F}$,
$\operatorname{Ker}(\widehat v)$ est isomorphe à
$\widehat{\operatorname{Ker}(v)}$ et comme le support de
$\operatorname{Ker}(v)$ est contenu dans celui de $\mathcal F$, on voit que
$\operatorname{Ker}(u)$ est algébrisable~; démonstration analogue pour
$\operatorname{Im}(u)$ et $\operatorname{Coker}(u)$.
\end{proof}

\begin{proposition}[5.2.3]
\label{III.5.2.3-fr}
Soient $A$ un anneau adique noethérien, $\mathfrak J$ un idéal de définition de
$A$, $\mathfrak Y=\operatorname{Spf}(A)$,
$f:\mathfrak X\to\mathfrak Y$ un morphisme propre de préschémas formels. On pose
$Y_k=\operatorname{Spec}(A/\mathfrak J^{k+1})$,
$X_k=\mathfrak X\times_{\mathfrak Y}Y_k$, et pour tout
$\mathcal O_{\mathfrak X}$-Module $\mathcal F$,
\[
\mathcal F_k=\mathcal F\otimes_{\mathcal O_{\mathfrak X}}\mathcal O_{X_k}
=\mathcal F/\mathfrak J^{k+1}\mathcal F.
\]
Soit $\mathcal L$ un $\mathcal O_{\mathfrak X}$-Module inversible, et supposons que
$\mathcal L_0=\mathcal L/\mathfrak J\mathcal L$ soit un
$\mathcal O_{X_0}$-Module ample~; pour tout $\mathcal O_{\mathfrak X}$-Module
$\mathcal F$ et tout entier $n$, posons
$\mathcal F(n)=\mathcal F\otimes\mathcal L^{\otimes n}$. Alors, pour tout
$\mathcal O_{\mathfrak X}$-Module cohérent $\mathcal F$, il existe un entier $n_0$
tel que, pour tout $n\geq n_0$, les propriétés suivantes aient lieu~:
\begin{enumerate}
\item[(i)] L'homomorphisme canonique
$\mathrm H^0(\mathfrak X,\mathcal F(n))\to
\mathrm H^0(\mathfrak X,\mathcal F_k(n))$ est surjectif pour tout $k\geq0$.
\item[(ii)] On a $\mathrm H^q(\mathfrak X,\mathcal F(n))=0$ pour tout $q>0$.
\end{enumerate}
\end{proposition}

\begin{proof}
On sait que les espaces sous-jacents à $\mathfrak X$ et à $X_0$ sont les mêmes~;
les faisceaux
\[
\mathcal M_k=\mathfrak J^k\mathcal F/\mathfrak J^{k+1}\mathcal F
\]
étant annulés par $\mathfrak J$, peuvent être considérés comme des
$\mathcal O_{X_0}$-Modules cohérents ($0_{\mathrm I}$, 5.3.10)~; en outre, si on
pose $\mathcal M_k(n)=\mathcal M_k\otimes_{\mathcal O_{X_0}}
\mathcal L_0^{\otimes n}$, on voit aussitôt que
\[
\mathcal M_k(n)=\mathfrak J^k\mathcal F(n)/
\mathfrak J^{k+1}\mathcal F(n).
\]
Notons que, puisque $\mathcal L_0$ est ample pour $f_0:X_0\to Y_0$, et

